\section{Постановка задачи оптимизации}
\label{sec:optimization_problem}

\subsection{Показатель покрытия}
\label{subsec:coverage_metric}

Рассматривается одноэшелонное Walker-построение круговых орбит с вектором
проектных переменных
\begin{equation}
    \mathbf{x}
    =\left(\mathbf{x}_{\mathrm{c}},\mathbf{x}_{\mathrm{d}}\right)
    =\left(h,i,\Delta\Omega,f,S,P\right),
    \label{eq:design_variables}
\end{equation}
где
\begin{equation}
    \mathbf{x}_{\mathrm{c}}
    =\left(h,i,\Delta\Omega,f\right)\in\mathcal{X}_{\mathrm{c}}
    \subset\mathbb{R}^{4},
    \qquad
    \mathbf{x}_{\mathrm{d}}
    =\left(S,P\right)\in\mathcal{X}_{\mathrm{d}}
    \subset\mathbb{N}_{++}^{2}.
    \label{eq:design_variable_domains}
\end{equation}
Здесь $h$ -- высота орбиты, $i$ -- наклонение, $\Delta\Omega$ -- шаг по
долготе восходящего узла, $f$ -- фазовый сдвиг между соседними плоскостями,
$S$ -- число КА в плоскости, а $P$ -- число плоскостей. Множества
$\mathcal{X}_{\mathrm{c}}$ и $\mathcal{X}_{\mathrm{d}}$ задают допустимые
значения соответственно непрерывных и целочисленных проектных переменных.

Пусть $\mathcal{R}$ -- заданная область земной поверхности. Используя
обозначения подраздела~\ref{subsec:geoprotection}, через
$\mathcal{C}(t,\mathbf{x})\subseteq\mathcal{R}$ обозначим множество точек
$Q$, для каждой из которых в момент $t$ существует хотя бы один КА $S_k$,
одновременно удовлетворяющий угловому критерию поля
зрения~\eqref{eq:angular_coverage_criterion}, векторному критерию прямой
видимости~\eqref{eq:vector_coverage_criterion} и, если учитывается защита
геостационарных систем, условию
геозащиты~\eqref{eq:geoprotection_condition}. Если геозащитная маска не
используется, последнее условие не применяется. Полуугол поля зрения
задаётся непосредственно либо определяется из условия
$\mathrm{SNR}\geq\mathrm{SNR}_{\min}$ согласно подразделу~\ref{subsec:snr}
и затем ограничивается эффективным углом $\alpha^*$
из~\eqref{eq:effective_fov}.

Обозначим через $A(\mathcal{D})$ площадь произвольной области
$\mathcal{D}$ на поверхности Земли. Тогда мгновенный показатель покрытия
определяется как
\begin{equation}
    \mu(t,\mathbf{x})
    =\frac{A\!\left(\mathcal{C}(t,\mathbf{x})\right)}{A(\mathcal{R})},
    \qquad \mu(t,\mathbf{x})\in[0,1].
    \label{eq:coverage_metric}
\end{equation}
Таким образом, $\mu(t,\mathbf{x})$ представляет долю площади заданной
области, покрытую группировкой в момент времени $t$.

\subsection{Формулировка задачи оптимизации}
\label{subsec:optimization_formulation}

Пусть $\mathcal{I}=[t_{\mathrm{start}},t_{\mathrm{end}}]$ -- непрерывный
интервал, на котором требуется обеспечить покрытие. Тогда задача оптимизации
формулируется следующим образом:
\begin{equation}
\label{eq:optimization_problem}
\begin{aligned}
    \min_{\mathbf{x}}\quad
        & J(\mathbf{x}) = SP,\\
    \text{s.t.}\quad
        & \inf_{t\in\mathcal{I}}\mu(t,\mathbf{x}) = 1,\\
        & r_{\mathrm{req}} - r_{\min}^{\mathrm{Walker}}(\mathbf{x}) \leq 0,\\
        & \mathbf{x}_{\mathrm{c}}\in\mathcal{X}_{\mathrm{c}},
        \qquad \mathbf{x}_{\mathrm{d}}\in\mathcal{X}_{\mathrm{d}},
\end{aligned}
\end{equation}
Таким образом, требуется минимизировать общее число КА при полном покрытии
области $\mathcal{R}$ на всём рассматриваемом интервале и при соблюдении
заданного минимального межспутникового расстояния.
